\documentclass[11pt, a4paper]{article}

% --- UNIVERSAL PREAMBLE BLOCK (MANDATORY FOR COMPILATION) ---
\usepackage[a4paper, top=2.5cm, bottom=2.5cm, left=2cm, right=2cm]{geometry}
\usepackage{fontspec}
\usepackage{amssymb} % Pentru simboluri precum \aleph, \mathfrak
\usepackage{graphicx}
\usepackage{hyperref}

% Babel setup for Romanian (main) and Greek
\usepackage[romanian, bidi=basic, provide=*]{babel}
% \babelprovide[import, onchar=ids fonts]{greek}
% \babelprovide[import, onchar=ids fonts]{english}

% Set default/Latin font to Sans Serif (Noto Sans)
\babelfont{rm}{Noto Sans}
% Assign Noto Sans for Greek as well to ensure consistent look
\babelfont[greek]{rm}{Noto Sans}

% Fix lists for non-English main language
\usepackage{enumitem}
\setlist[itemize]{label=-}

% Definire comandă pentru compatibilitate cu sintaxa ta \textgreek
\newcommand{\textgreek}[1]{\foreignlanguage{greek}{#1}}

% --- END PREAMBLE ---

\title{Analiză Textuală Critică: Ioan 18:33-38}
\author{Țacu Darius-Ștefan \\ Grupa 7, An I Teologie Pastorală}
\date{}

\begin{document}

\maketitle

\section{Analiză pe Versete}

\subsection{Ioan 18:33}
\textbf{Text grec:} \\
\textgreek{Εἰσῆλθεν οὖν ⸂πάλιν εἰς τὸ πραιτώριον⸃ ▫ὁ Πιλᾶτος καὶ ἐφώνησε τὸν Ἰησοῦν καὶ εἶπεν αὐτῷ· σὺ εἶ ὁ βασιλεὺς τῶν Ἰουδαίων;}

\textbf{Traducere:} \\
\textit{Deci a intrat iarăși în pretoriu Pilat și a chemat pe Iisus și I-a zis: Tu ești împăratul Iudeilor?}

\begin{itemize}
    \item Semnul \textgreek{⸂} încadrează expresia \textgreek{πάλιν εἰς τὸ πραιτώριον}. Aparatul critic notează o inversiune (\textit{în pretoriu iarăși}) în anumite manuscrise.

    \item Semnul \textgreek{▫} înainte de \textgreek{ὁ Πιλᾶτος} arată ca unele manuscrise, inclusiv $\mathfrak{P}^{66}$, omit articolul hotărât \textgreek{ὁ} înainte de numele propriu.
\end{itemize}

\subsection{Ioan 18:34}
\textbf{Text grec:} \\
\textgreek{ἀπεκρίθη ⸆ Ἰησοῦς· ἀπὸ σεαυτοῦ σὺ τοῦτο λέγεις ἢ ἄλλοι εἶπόν σοι περὶ ἐμοῦ;}

\textbf{Traducere:} \\
\textit{Răspuns-a Iisus: De la tine însuți zici aceasta, sau alții ți-au spus-o despre Mine?}

\begin{itemize}
    \item \textbf{Adăugire:} Semnul \textgreek{⸆} indică o posibilă inserție după verbul \textgreek{ἀπεκρίθη}.  Aparatul critic arată varianta \textgreek{αυτω} (lui). Multe manuscrise (Textul Majoritar $\mathfrak{M}$, $\aleph$, A) adaugă pronumele: \textit{i-a răspuns lui}.
\end{itemize}

\subsection{Ioan 18:35}
\textbf{Text grec:} \\
\textgreek{ἀπεκρίθη ὁ Πιλᾶτος· μήτι ἐγὼ Ἰουδαῖός εἰμι; τὸ ἔθνος τὸ σὸν καὶ ⸂οἱ ἀρχιερεῖς⸃ ⸀παρέδωκάν σε ἐμοί· τί ἐποίησας;}

\textbf{Traducere:} \\
\textit{Răspuns-a Pilat: Nu cumva sunt și eu iudeu? Neamul Tău și arhiereii Te-au predat mie. Ce ai făcut?}

\begin{itemize}
    \item Semnul \textgreek{⸀} la \textgreek{παρέδωκάν} (plural): unele manuscrise (ex. $\aleph$*) au varianta la singular \textgreek{παρέδωκεν} (a predat)
\end{itemize}

\subsection{Ioan 18:36}
\textbf{Text grec:} \\
\textgreek{ἀπεκρίθη Ἰησοῦς· ἡ ⸂βασιλεία ἡ ἐμὴ⸃ οὐκ ἔστιν ἐκ τοῦ κόσμου τούτου· εἰ ἐκ τοῦ κόσμου τούτου ἦν ἡ βασιλεία ἡ ἐμή, οἱ ὑπηρέται οἱ ἐμοὶ ἠγωνίζοντο [ἂν] ἵνα μὴ παραδοθῶ τοῖς Ἰουδαίοις· νῦν δὲ ἡ βασιλεία ἡ ἐμὴ οὐκ ἔστιν ἐντεῦθεν.}

\textbf{Traducere:} \\
\textit{Iisus a răspuns: Împărăția Mea nu este din lumea aceasta. Dacă Împărăția Mea ar fi din lumea aceasta, slujitorii Mei s-ar fi luptat ca să nu fiu predat iudeilor. Dar acum Împărăția Mea nu este de aici.}

\begin{itemize}
    \item Semnul \textgreek{⸂} la \textgreek{βασιλεία ἡ ἐμὴ} (Împărăția Mea). Codex Sinaiticus ($\aleph$) are ordinea inversă: \textgreek{εμη βασιλεια}.
\end{itemize}

\subsection{Ioan 18:37-38}
\textbf{Text grec (v.37):} \\
\textgreek{εἶπεν οὖν αὐτῷ ὁ Πιλᾶτος· οὐκοῦν βασιλεὺς εἶ σύ; ἀπεκρίθη ὁ Ἰησοῦς· σὺ λέγεις ὅτι βασιλεύς εἰμι. ἐγὼ εἰς τοῦτο γεγέννημαι καὶ εἰς τοῦτο ἐλήλυθα εἰς τὸν κόσμον ἵνα μαρτυρήσω τῇ ἀληθείᾳ· πᾶς ὁ ὢν ἐκ τῆς ἀληθείας ἀκούει μου τῆς φωνῆς.}

\textbf{Text grec (v.38):} \\
\textgreek{λέγει ⸆ αὐτῷ ▫ὁ Πιλᾶτος· τί ἐστιν ἀλήθεια; Καὶ τοῦτο εἰπὼν πάλιν ἐξῆλθεν πρὸς τοὺς Ἰουδαίους καὶ λέγει αὐτοῖς· ἐγὼ οὐδεμίαν εὑρίσκω ἐν αὐτῷ αἰτίαν.}

\textbf{Traducere (v.37-38):} \\
\textit{Deci i-a zis Pilat: Așadar ești Împărat? Răspuns-a Iisus: Tu zici că Eu sunt Împărat. Eu spre aceasta M-am născut și pentru aceasta am venit în lume, ca să mărturisesc adevărul; oricine este din adevăr ascultă glasul Meu. \\
    I-a zis Pilat: Ce este adevărul? Și zicând aceasta, a ieșit iarăși la iudei și le-a zis: Eu nu găsesc în El nicio vină.}

\begin{itemize}
    \item \textbf{v.37:} Variații minore legate de particule și poziția pronumelui \textgreek{ἐγὼ} ("eu").
    \item Semnul \textgreek{▫} indică lipsa articolului \textgreek{ὁ} în manuscrise importante precum B.
\end{itemize}
\end{document}